% Part: second-order-logic
% Chapter: metatheory
% Section: undecidability-and-axiomatizability

\documentclass[../../../include/open-logic-section]{subfiles}

\begin{document}

\olfileid{sol}{met}{nax}

\olsection{二阶逻辑不可公理化}

\begin{thm}
\ollabel{thm:sol-undecidable}
二阶逻辑是不可判定的。
\end{thm}

\begin{proof}
一阶语句在一阶逻辑中有效当且仅当它在二阶逻辑中有效，而一阶逻辑是不可判定的。
\end{proof}

\begin{thm}
\ollabel{cor:sol-not-axiomatizable}
不存在二阶逻辑的可靠且完备的推导系统。
\end{thm}

\begin{proof}
设 $!A$ 为算术语言中的一个语句。$\Sat{N}{!A}$ 当且仅当 $\Th{PA^2} \Entails !A$。令 $!P$ 为 $\Th{PA^2}$ 的九条公理的合取。$\Th{PA^2} \Entails !A$ 当且仅当 $\Entails !P \lif !A$，即对任意结构 $M$，$\Sat{M}{!P \lif !A$}。现在考虑由原语句经如下替换并加以全称量化得到的语句 $\lforall[z][\lforall[u][\lforall[u'][\lforall[u''][\lforall[L][(!P' \lif !A')]]]]]$：将 $\Obj{0}$ 替换为 $z$，$\prime$ 替换为一元函数变元 $u$，$+$ 与 $\times$ 分别替换为二元函数变元 $u'$ 与 $u''$，$<$ 替换为二元关系变元 $L$，然后对这些替换引入的变元进行全称量化。该语句是纯二阶逻辑中的有效语句，当且仅当原语句有效，当且仅当 $\Th{PA^2} \Entails !A$，当且仅当 $\Sat{N}{!A}$。因此，如果存在二阶逻辑的可靠且完备的证明系统，我们就可用它来定义 $\Struct{N}$ 中真语句的一个可计算枚举 $f\colon \Nat \to \Sent[L_A]$。该函数在 $\Th{Q}$ 中可由某个一阶公式 $!B_f(x, y)$ 表示。于是，公式 $\lexists[x][!B_f(x, y)]$ 将定义 $\Struct{N}$ 中一阶真语句的集合，这与 塔斯基 定理矛盾。
\end{proof}



\end{document}